\section{Experiments}

\subsection{Training Procedure}
\todo{Summarise shared training setup (optimiser, epochs, early stopping, same conditions across models).}

\subsubsection*{Pinball vs Revenue Weighted Pinball Loss}

Quantile regression is trained using the pinball loss, which estimates conditional quantiles \( \tau \in (0,1) \):

\begin{equation}
\mathcal{L}_{\tau}(\hat{y}, y) =
\begin{cases}
\tau (y - \hat{y}) & \text{if } y \geq \hat{y} \\
(1 - \tau)(\hat{y} - y) & \text{otherwise}
\end{cases}
\end{equation}

This asymmetric loss enables direct learning of prediction intervals without assuming a parametric distribution.

However, standard pinball loss treats all items equally, which is unsuitable for retail forecasting. Forecast errors have unequal business costs: underestimating high-revenue products leads to unmet demand and lost sales, while overestimation leads to excess inventory. These costs scale with price and volume.

We therefore use a revenue-weighted variant:

\begin{equation}
\mathcal{L}_{\text{weighted}} =
\sum_{i=1}^{N} w_i \sum_{\tau \in \mathcal{T}} \mathcal{L}_{\tau}(\hat{y}_{i,\tau}, y_i),
\quad \sum_i w_i = 1
\end{equation}

where \( w_i \propto \) (sales volume \( \times \) average price).

\paragraph{} Equal weighting optimises average error; revenue weighting optimises financial impact.

\subsection{Evaluation Metrics}
\todo{Define metrics (RMSE, MAE, weighted variants, pinball, coverage, interval width).}